% Options for packages loaded elsewhere
\PassOptionsToPackage{unicode}{hyperref}
\PassOptionsToPackage{hyphens}{url}
%
\documentclass[
]{article}
\usepackage{amsmath,amssymb}
\usepackage{lmodern}
\usepackage{iftex}
\ifPDFTeX
  \usepackage[T1]{fontenc}
  \usepackage[utf8]{inputenc}
  \usepackage{textcomp} % provide euro and other symbols
\else % if luatex or xetex
  \usepackage{unicode-math}
  \defaultfontfeatures{Scale=MatchLowercase}
  \defaultfontfeatures[\rmfamily]{Ligatures=TeX,Scale=1}
\fi
% Use upquote if available, for straight quotes in verbatim environments
\IfFileExists{upquote.sty}{\usepackage{upquote}}{}
\IfFileExists{microtype.sty}{% use microtype if available
  \usepackage[]{microtype}
  \UseMicrotypeSet[protrusion]{basicmath} % disable protrusion for tt fonts
}{}
\makeatletter
\@ifundefined{KOMAClassName}{% if non-KOMA class
  \IfFileExists{parskip.sty}{%
    \usepackage{parskip}
  }{% else
    \setlength{\parindent}{0pt}
    \setlength{\parskip}{6pt plus 2pt minus 1pt}}
}{% if KOMA class
  \KOMAoptions{parskip=half}}
\makeatother
\usepackage{xcolor}
\setlength{\emergencystretch}{3em} % prevent overfull lines
\providecommand{\tightlist}{%
  \setlength{\itemsep}{0pt}\setlength{\parskip}{0pt}}
\setcounter{secnumdepth}{-\maxdimen} % remove section numbering
\ifLuaTeX
  \usepackage{selnolig}  % disable illegal ligatures
\fi
\IfFileExists{bookmark.sty}{\usepackage{bookmark}}{\usepackage{hyperref}}
\IfFileExists{xurl.sty}{\usepackage{xurl}}{} % add URL line breaks if available
\urlstyle{same} % disable monospaced font for URLs
\hypersetup{
  hidelinks,
  pdfcreator={LaTeX via pandoc}}

\author{Dietmar Gerald Schrausser}
\date{09.05.2026}


\begin{document}

\hypertarget{foliumblat-vir}{%
\section{Foliu{[}m{]}/Blat VIr}\label{foliumblat-vir}}

A \emph{contemporary} English translation is presented from comparing
the existing Latin, Early New High German (ENHG) and English
(translation) texts for better comprehensibility. Especially since the
ENHG text is difficult to understand compared to modern High German, as
a \emph{profound} knowledge of German (as a mother tongue) as well as
already aknowledge of various \emph{dialects} is required.

\hypertarget{certain-things-need-to-be-explained-about-the-planetary-sphere-and-the-orbits.}{%
\subsection{Certain things need to be explained about the planetary
sphere and the
orbits.}\label{certain-things-need-to-be-explained-about-the-planetary-sphere-and-the-orbits.}}

\hypertarget{the-difference-between-celestial-and-elementary-orbits1}{%
\subsubsection[The difference between celestial and elementary
orbits]{\texorpdfstring{The difference between celestial and elementary
orbits\footnote{`orbitu{[}m{]}' and `umbkreys'.}}{The difference between celestial and elementary orbits}}\label{the-difference-between-celestial-and-elementary-orbits1}}

The entire mechanism of the \emph{corpus}\footnote{`COrporalis' and
  `leiplich'.} of the world consists of \emph{two} things: namely, (I)
\emph{celestial} nature and (II) \emph{elemental} nature. Furthermore,
the heaven (I) is divided into \emph{three} \emph{principal heavens},
namely the (i) \emph{fiery}\footnote{`empyreum' and `feürigen'.}, the
(ii) \emph{crystalline}, and the (iii) \emph{firmament}.

Within the firmament (iii), which \emph{indeed}\footnote{`vero'.} is the
heaven of the stars\footnote{`celu{[}m{]} stellarum' and `gestirnt
  himel'.}, are the seven orbits of the seven planets: Saturn, Jupiter,
Mars, Sun, Venus, Mercury and Moon. However, the first part of
primordial matter\footnote{`materie prime' and `erste{[}n{]} materi'},
which according to the philosopher\footnote{`s{[}ecundum{]}
  philosophu{[}m{]}' and `nach sag des weysen', in this context, it
  almost always refers to Aristotle.} is divided into two orbits, is
called the crystalline or aqueous\footnote{`aquei', in medieval thought,
  this was the sphere above the stars, often associated with the
  \emph{waters above the firmament} mentioned in Genesis.} heaven (ii).
The upper orbit is the \emph{primum mobile}. The nature of this orbit is
that all of them are in motion\footnote{`mouent{[}ur{]}' and `bewegt
  werden'.}, except for the fiery heaven (i), which is at rest.

In fact\footnote{`vero' and `Aber'.}, the \emph{elemental} nature (II)
can be divided into \emph{four} main spheres: (i) fire, (ii) air, (iii)
earth, and (iv) water. The sphere of fire (i) has three sections: (a)
the \emph{uppermost}, which is called \emph{fiery}, and the
\emph{middle} and \emph{lowermost}, which is called Olympian\footnote{`olimpium'
  and `liecht'.}.

Similarly, the air (ii) has three sections, the (a) \emph{uppermost},
which is called \emph{ethereal}\footnote{`ethereum' and `scheynlich'.},
the (b) \emph{middle}, and the (c) \emph{lowest}, which
is/are\footnote{`dicit{[}ur{]}', singular vs.~`der mittel vn{[}d{]}
  underst lüftig'.} called \emph{airy}\footnote{`aereum' and `lüftig'.}.
In the uppermost area (a) there is warmth and light due to the proximity
of the sun, a similar situation exists in the lowermost part (c), but
this is due to the \emph{reflection}\footnote{`rep{[}er{]}cussionem' and
  `wider scheyns'.} of the \emph{rays from the earth}\footnote{`ratio{[}rum{]}
  \ldots{} a terra' and `glentz von der erden'.}. However, in the middle
section (b), which the \emph{reflected rays}\footnote{`rep{[}er{]}cussio
  radio{[}rum{]}' and `widerscheyn d{[}er{]} glentz'.} cannot reach,
there is cold and darkness. It is said that this is inhabited by
\emph{demons}\footnote{`demones' and `teüfel'.} who have been cast
into\footnote{`detrusi' and `verstössen'.} this \emph{dark}\footnote{`caliginosum'
  and `tunckeln'.} \emph{air}.\footnote{found particularly in the
  handbook of Christian doctrine, the `Elucidarium', a frequent
  theological reference to the \emph{atmospheric prison} in which fallen
  angels remain until the Last Judgment, c.f. Ephesians or 2 Peter 2:4
  `for God did not spare the angels who sinned, but cast them into
  \emph{Hell} and put them in \emph{chains of darkness}, {[}\ldots{]}'.
  Note especially the `\emph{chains of darkness},' which symbolize
  complete captivity and \emph{separation from} \emph{God's light}.}
Severe weather events\footnote{`tempestates' and `ungestümigkeit'.} such
as thunder, hail, snow and similar phenomena also occur
there.\footnote{c.f. Francesco Maria Guazzo
  (\href{https://archive.org/details/compendiummalefi00guaz/page/n2/mode/1up}{1608}).}

From these, one forms \emph{twelve orbits}\footnote{`orbes' and
  `umbkrais'.} of the world that surround \emph{earth and water} (iii,
iv), which can all be called \emph{heavens}. But these are
surpassed\footnote{`excellit' and `ubertrifft'.} by the \emph{heaven of
the Trinity}, \emph{God Himself}\footnote{`ipse deus' and `der got'.},
who is \emph{among all} things\footnote{`in omnibus' and `in allen'.}
and \emph{beyond all} things\footnote{`super om{[}n{]}ia' and `über
  alle'.}.

The distances of the \emph{forementioned}\footnote{`predictorum' and
  `vorgenante{[}n{]}'.} orbits and planets from the Earth to the Moon
are 15625 miles; \emph{miles, that is} \textbf{\emph{126}}
\emph{stadia}\footnote{`miliaria. hec sunt stadia .cxxvi.', this results
  in 15625 miles × 126 = 1968750 stadia, which, using \emph{Attic}
  stadia, gives a distance of 0.18498 × 1968750 = 364179.375 km from the
  Earth to the Moon. With a current \emph{maximum} distance
  (\emph{apogee}) of 406740 km, the difference is -42560.625 km, and
  with a \emph{minimal} distance (\emph{perigee}) of 356410 km, the
  difference is +7769.375 km. However, this corresponds to the distance
  of the moon about \emph{800 million} years ago (\emph{Tonium} in the
  \emph{Neoproterozoic} era, 363000 to 367000 km).}. The distance from
the Moon to Mercury is 7812.5 miles. The distance from Mercury to Venus
is the same. The distance from Venus to the Sun is 22436 miles. From the
Sun to Mars it is 15625 miles. From Mars to Jupiter it is 7812 miles.
The distance from Jupiter to Saturn is the same. From Saturn to the
firmament it is 22436 miles. Therefore, the distance from Earth to the
stars is 109375 miles.

\hypertarget{the-distinction-between-the-heavenly-hierarchies-power-or-principalities}{%
\subsubsection{The distinction between the heavenly hierarchies {[}power
or
principalities{]}}\label{the-distinction-between-the-heavenly-hierarchies-power-or-principalities}}

Regarding the \emph{heavenly} nature (I), several authors have indeed
made a \emph{threefold} distinction between a (i) \emph{superheavenly}
(or the \emph{supernatural}), a (ii) \emph{heavenly} and a (iii)
\emph{subheavenly}.

The superheavenly hierarchy (i) exists in three persons, as some have
claimed, \emph{however erroneously}, because according to
\textbf{Dionysius}, \emph{hierarchy} implies an order, but such an order
does not exist absolutely in the three persons, except as a natural
order.\footnote{`\ldots{} vt q{[}ui{]}dam dixerunt {[}et{]} male
  \ldots{}', Pseudo-Dionysius the Areopagite
  (\href{https://digital.bodleian.ox.ac.uk/objects/ee19a692-5066-4c84-95db-d93c4fc82b97/}{1350})
  defined whether the term hierarchy can be applied to the Holy Trinity.}

The \emph{heavenly} (ii) is divided into the \emph{Order of Angels}, the
\emph{subheavenly} (iii) into \emph{holy persons}. Furthermore, the
\emph{heavenly} hierarchy (ii) is divided into the (a) upper, the (b)
middle and the (c) lower levels. The upper level (a) comprises three
orders, the (a1) \emph{Seraphim}, the (a2) \emph{Cherubim} and the (a3)
\hspace{0pt}\hspace{0pt}\emph{Throne Angels}.

The first (a1) consider\footnote{`{[}con{]}siderant' and `betrachten'.}
God's \emph{benevolence}\footnote{`bonitate{[}m{]}' and `guttheit'.},
the second (a2) his \emph{greatness}\footnote{`virtute{[}m{]}' and
  `kraft'.}, the third (a3) \hspace{0pt}\hspace{0pt}his
\emph{appropriateness}\footnote{`equitate{[}m{]}' and `gleicheit'.}.
Likewise, God \emph{loves} in the first (a1) through
\emph{charity}\footnote{`charitas' and `lib'.}, in the second (a2) he
understands through \emph{truth}, in the third (a3)
\hspace{0pt}\hspace{0pt}he resides\footnote{`sedet' and `sitzt'.} in
\emph{justice}\footnote{`equitas' and `gleicheit'.}.\footnote{by
  Pseudo-Dionysius the Areopagite
  (\href{https://digital.bodleian.ox.ac.uk/objects/ee19a692-5066-4c84-95db-d93c4fc82b97/}{1350}),
  often also attributed to Saint Gregory the Great (c.f. Étaix,
  \href{https://books.google.com/books?id=r4ERnwEACAAJ}{1999}), also
  adapted by Meister Eckhart
  (\href{http://www.eckhart.de/opus.htm}{1936}), reflects his
  Trinitarian theology, in which he harmonizes the different ways of
  \emph{perceiving God} with the three \emph{persons} of the Trinity.}

The middle hierarchy (b) comprises (b1) \emph{Ruler-}\footnote{`dominat{[}i{]}ones'
  and `herschengel'.}, (b2) \emph{Prince-}\footnote{`principatus' and
  `fürstengel'.} and (b3) \emph{Power Angels}\footnote{`potestates' and
  `gewaltengel', better \emph{ability}, \emph{skill}.}. The first (b1)
govern the tasks of the angels, the \emph{following}\footnote{`Sequentes'
  and `die andern'.} are \emph{Advisors}\footnote{`capitibunt
  p{[}re{]}sunt' and `pflegen der öbern'.} to the leaders of the
peoples. The \emph{last ones} (b3) suppress\footnote{`coercent' and
  `zwinge{[}n{]}'.} the \emph{ability}\footnote{`potestatem' and
  `macht'.} of the \emph{`devils'}\footnote{`demonu{[}m{]}' and
  `teüfel', \emph{the criminals}.}. Likewise, the Lord \emph{rules}
among the first ones (b1) as \emph{majesty}, among the
following\footnote{`in sequentibus' and `in den andern'.} (b2) he
\emph{rules} through a \emph{principality}\footnote{`principat{[}us{]}'
  and `fürstenthumb'.} and by the last ones (b3) his \emph{`health' is
kept}\footnote{`tenet{[}ur{]} vt salus' and `gehalte{[}n{]} als das
  hail'.}.

The lower hierarchy (c) also includes three orders, (c1)
\emph{Virtues}\footnote{`virtutes' and `kreftengel'.}, (c2)
\emph{Archangels} and (c3) \emph{Angels}. The Virtues (c1) are
responsible for performing great miracles, the archangels (c2) for
\emph{proclaiming} important matters, and the angels (c3) for the
\emph{protection} of people. Thus, in the first (c1) God acts as a
\emph{force}\footnote{`virt{[}us{]}' and `kraft'.}, in the
\emph{following} (c2) he reveals himself as \emph{light}, in the third
(c3) he \emph{nourishes} as \emph{inspiration}\footnote{`inspira{[}n{]}s'
  and `eyngeystender'.}.

This according to {[}Pope{]} \textbf{Gregor}.\footnote{`Hec dicta sunt
  s{[}ecundu{]}m Gregoriu{[}m{]}', common phrase in medieval
  manuscripts, the preceding passage was marked as quoted from the work
  of Pope Gregory I.} According to \textbf{Dionysius}, however, the
\emph{Virtues} (c1) form the middle order of the second hierarchy (b2),
while the \emph{Prince Angels} (b2) form the first order of the third
hierarchy (c1). This must be taken into account so that the
\emph{Trinity of the divine Person}\footnote{`trinitas p{[}er{]}sonarum
  diuinarum' and `trinitet der gottliche{[}n{]} person'.} becomes
visible in \emph{each} of the \emph{three} hierarchies mentioned (i, ii,
{[}iii{]}), on all levels, the upper (1), middle (2) and lower (3)
equally, as\footnote{`vt patet'.} in the heavenly hierarchy
(ii).\footnote{Thomas Aquinas
  (\href{https://doi.org/10.3931/e-rara-15170}{1485}) elaborated on this
  in works such as the `Summa Theologica' and followed the Dionysian
  structure. The idea that every hierarchy reflects the \emph{Trinity of
  divine persons} is a central theme of medieval Scholasticism.}

\hypertarget{from-time-or-ages}{%
\subsubsection{From time or ages}\label{from-time-or-ages}}

The ages of the world are understood symbolically according to the
\emph{ages of human life}. Of the six ages of the world\footnote{The
  Christian understanding of history popularized by Saint Augustine
  (\href{https://doi.org/10.3931/e-rara-12425}{1489}). In `De civitate
  Dei', he laid the foundation for this chronological framework during
  the decline of the Western Roman Empire as a comprehensive Christian
  interpretation of \emph{human history} and the relationship between
  the \emph{spiritual and secular} worlds.}, the first (I) begins with
the \emph{creation of the world} and lasts until the Great Flood.
According to the Hebrew account\footnote{`hebraicam veritate{[}m{]}' und
  `nach hebreyscher warheit', referred to the \emph{original} Hebrew
  text of the Bible, as opposed to the Greek \emph{Septuagint}.}, it was
1656 years\footnote{`.i656.' vs.~`.im.vc. lvi.', 1656 vs.~1556.}, but
according to the \textbf{Septuagint} (as \textbf{Isidore} explains) and
many others whom we follow here, it was 2242 years.

And so they differ by 586 years, which the Hebrews have fewer in this
age.\footnote{Characteristic of Isidore of Seville
  (\href{https://books.google.com/books?id=wb8Z_vlSyJwC}{1802}) or Bede
  (\href{http://archivesetmanuscrits.bnf.fr/ark:/12148/cc76825p}{1125}),
  they harmonized the different time periods.} According to this
calculation, \emph{Methuselah}\footnote{`matusale mortuus' and
  `matusaleh gestorben', a common chronological problem in medieval
  research: whether Methuselah survived the Flood.} died before the
Flood, but in the same year in which the Flood took place.

The second\footnote{`S{[}e{]}c{[}un{]}da' and `Das andre'.} age (II)
begins with the Flood and lasts until the birth of Abraham, according to
the Hebrews it lasted 292 years\footnote{`292' vs.~`.ijc.lxxxij.', 292
  vs. 282.}, according to the Septuagint 942 years. Thus, the difference
is 650 {[}\emph{660}{]} years, which the Hebrews have in fewer years.
But I \emph{could not find out}\footnote{`inuenire no{[}n{]} potui' and
  `nicht möge{[}n{]} finden'.} what the reason for these large
differences is.

The Third Age (III) begins with the birth of Abraham and lasted until
the beginning of David's reign. According to the Hebrews, it comprised
941 years; according to the Septuagint, 940 years. The fourth age (IV)
begins with the reign of David and lasts until the Babylonian
captivity\footnote{`transmigrat{[}i{]}o{[}n{]}is babilonis' and
  `übergang babilonis'.}, it comprises 484 years according to the
Hebrews and 485 years according to the Septuagint.

The fifth age (V) begins with the Babylonian captivity, more precisely
with the destruction of Jerusalem and the destruction of the temple
within it, and lasted until the birth of Christ. According to the
aforementioned mode, which we are following here, it comprises 590
years. There is considerable disagreement regarding the calculation of
the years of this era; some calculate them in different ways.

The sixth age (VI) begins with the birth of Christ and lasts until the
\emph{end of the world}, the \emph{date}\footnote{`terminum' and `zil'.}
of which only God knows. This is called \emph{old age}\footnote{`senectus'
  and `das alt'.} or \emph{the last hour}\footnote{`hora nouissima' and
  `letzt stu{[}n{]}d', known as the first sentence in `De Contemptu
  Mundi' of Bernard of Cluny
  (\href{https://books.google.com/books?id=pAM_lfiF7EUC}{1864}), in
  which he describes the corruption of the world and the approach of the
  end.}. To these ages, however, one may add a \emph{7th age}, that of
those who now \emph{rest}\footnote{`quiescentium' and `die nw ruen'.},
which takes place at the same time as the sixth age (VI). As well as an
\emph{eighth age}, that of those who \emph{resurrect}\footnote{`resurgentium'
  and `auffersteenden'.}.

Furthermore, according to the Hebrews, there are 10 generations in the
first age of the world (I), 10 in the second (II), 14 in the third
(III), and 17 in the fourth (IV). However, \textbf{Matthew} gives 14 to
emphasize the \emph{mystical significance}\footnote{`gratia misterij'
  and `aus verborgener bedeütnus', while historical records list 17
  generations in the fourth age, the Gospel of Matthew (c f.~Jiménez de
  Cisneros, \href{https://doi.org/10.3931/e-rara-46695}{1517}, Matthew
  1:17) deliberately divides Jesus' genealogy into three groups of 14
  each. This serves less as a \emph{historically accurate}
  representation than as a \emph{mystical} symbolization within the
  \emph{hermeneutic} technique of \emph{gematria} (where letters of the
  alphabet are assigned numerical values). It is noteworthy, however,
  that Matthew apparently \emph{omitted} several names from the
  historical records in order to \emph{be able} to divide the count into
  three groups of 14 each.}, and 14 in the fifth.

These are the ages of man: The (I) first is \emph{early
childhood}\footnote{`infantia' and `ungesprechheyt'.}, from birth to the
age of 7. The (II) second is \emph{childhood}\footnote{`puericia' and
  `kintheit'.} up to the age of 14. The (III) third is
\emph{`youth'}\footnote{`adolesce{[}n{]}tia' and `zeittigkeit'.}, from
the age of 15 to 38. The (IV) fourth is the \emph{young age}\footnote{`iuue{[}n{]}t{[}us{]}'
  and `iuge{[}n{]}t'.} up to the 49th year of life. The fifth (V) is
\emph{age}\footnote{`senect{[}us{]}' and `altheit'.}, from the 50th to
the 79th year of life. The (VI) sixth is \emph{old age}\footnote{`decrepita'
  and `das verlebt'.}, from the 80th year of life until the end of life.

\hypertarget{references}{%
\subsection{References}\label{references}}

Augustinus, A. (1489). \emph{De Civitate Dei}. Edited by Brant, S.,
Thomas, Nicolaus. Basel: Johannes Amerbach.
\url{https://doi.org/10.3931/e-rara-12425}

Beda. (1125). \emph{Beda, de Natura Rerum, de Temporibus, etc}. Latin
16361. Paris: Bibliothèque nationale de France. Département des
Manuscrits. \url{http://archivesetmanuscrits.bnf.fr/ark:/12148/cc76825p}

Bernard of Cluny. (1864). \emph{De Contemptu Mundi}. Harvard University:
J. T. Hayes. \url{https://books.google.com/books?id=pAM_lfiF7EUC}

De Seville, I. (1802). \emph{S. Isidori Hispalensis Episcopi Hispaniarum
Doctoris Opera Omnia}. Edited by Arevalo, F. Romae: Apvd Antonivm
Fvlgonivm. \url{https://books.google.com/books?id=wb8Z_vlSyJwC}

Dionysius. (1350). \emph{De Caelesti Hierarchia}. MS Gr 2. Oxford,
Magdalen College: Oxford Digital Library.
\url{https://digital.bodleian.ox.ac.uk/objects/ee19a692-5066-4c84-95db-d93c4fc82b97/}

Étaix, R. (1999). \emph{Gregorius Magnus, Homiliae in Evangelia}. Corpus
Christianorum Series Latina. Vol. 141. Turnhout: Brepols.
\url{https://books.google.com/books?id=r4ERnwEACAAJ}

Guazzo, F. M. (1608). \emph{Compendium Maleficarum in Tres Libros
Distinctum Ex Pluribus Authoribus}. Mediolani: Apud Haeredes August.
Tradati.
\url{https://archive.org/details/compendiummalefi00guaz/page/n2/mode/1up}

Jiménez de Cisneros, F. (1517). \emph{Biblia Polyglotta Complutensis}.
Complutum: Arnaldo Guillén de Brocar.
\url{https://doi.org/10.3931/e-rara-46695}

Meister Eckhart. (1936). \emph{Die Lateinischen Werke}. Edited by Weiss
Sturlese L. Stuttgart: Kohlhammer. \url{http://www.eckhart.de/opus.htm}

Thomas. (1485). \emph{Summa Theologica: Partes 1-3}. Basel: Michael
Wenssler. \url{https://doi.org/10.3931/e-rara-15170}

\end{document}
